\section*{2. “几何”情形}

我们将证明猜想 1.9 在等特征情形中成立. 关键的几何结果如下. 与 B. Saint-Donat 的交谈使我想到, 这一结果可能有用.

\subsubsection*{命题 2.1.}
\underline{设 $S$ 为代数闭域 $k$ 上的光滑连通射影曲线, $X$ 为 $k$ 上纯维数 $n+1$ 的光滑固有簇, $f:X\to S$ 为在 $X$ 的某个有限点集 $F$ 之外光滑的态射. 则 $S$, $X$ 以及 $f$ 的几何一般纤维 $X_{\bar\eta}$ 的平展上同调 Euler--Poincar\'e 示性数满足}
\begin{equation}\tag{2.1.1}
\chi(X)=\chi(S)\chi(X_{\bar\eta})
-(-1)^n\sum_{x\in F}\mu(X/S,x).
\end{equation}

$X$ 的 Euler--Poincar\'e 示性数是切丛的最高 Chern 类的迹（或在 $X$ 上的积分）:
\[
\chi(X)=\int_Xc_{n+1}(\Theta_X)
=(-1)^{n+1}\int_Xc_{n+1}(\Omega_X^1).
\]
记
\[
\Omega_X^1=\left(\Omega_X^1\otimes
(f^*\Omega_S^1)^{\otimes-1}\right)\otimes f^*\Omega_S^1.
\]
可逆层 $f^*\Omega_S^1$ 的全 Chern 类是 $S$ 上 $\Omega_S^1$ 的全 Chern 类的逆像:
\[
c^*(f^*\Omega_S^1)=1+f^*c^1(S,\Omega_S^1),
\]
而 $c^1$ 的杯积幂为零:
\[
\bigl(f^*c^1(S,\Omega_S^1)\bigr)^n
=f^*\bigl(c^1(S,\Omega_S^1)^n\bigr)=0
\qquad(n>1).
\]
因此
\begin{align*}
c^{n+1}(\Omega_X^1)
={}&c^n\!\left(\Omega_X^1\otimes
(f^*\Omega_S^1)^{\otimes-1}\right)
\smile f^*c_1(S,\Omega_S^1)\\
&+c^{n+1}\!\left(\Omega_X^1\otimes
(f^*\Omega_S^1)^{\otimes-1}\right).
\end{align*}

若 $S$ 的亏格为 $g$, 则 $c_1(S,\Omega_S^1)$ 等于 $S$ 的任一闭点 $s$ 的上同调类的 $(2g-2)$ 倍. 因此, 当 $X_s$ 非奇异时,
\begin{align*}
\mathrm I
&=\int_Xc^n\!\left(\Omega_X^1\otimes
(f^*\Omega_S^1)^{\otimes-1}\right)
\smile f^*c_1(S,\Omega_S^1)\\
&=(2g-2)\int_Xc^n\!\left(\Omega_X^1\otimes
(f^*\Omega_S^1)^{\otimes-1}\right)
\smile f^*\operatorname{cl}(s)\\
&=(2g-2)\int_{X_s}c^n\!\left(
X_s,\Omega_X^1\otimes(f^*\Omega_S^1)^{\otimes-1}
\otimes\mathcal O_{X_s}\right).
\end{align*}
此外, 在 $X_s$ 上, 层 $f^*\Omega_S^1$ 同构于 $\mathcal O_{X_s}$, 而 $\Omega_X^1\otimes_{\mathcal O_X}\mathcal O_{X_s}$ 出现在正合列
\[
0\longrightarrow f^*\Omega_S^1\otimes\mathcal O_{X_s}
\longrightarrow\Omega_X^1\otimes\mathcal O_{X_s}
\longrightarrow\Omega_{X_s}^1\longrightarrow0.
\]
中. 因此
\[
c^n\!\left(X_s,\Omega_X^1\otimes
(f^*\Omega_S^1)^{\otimes-1}\otimes\mathcal O_{X_s}\right)
=c^n(X_s,\Omega_{X_s}^1),
\]
及
\[
\mathrm I=(2g-2)\int_{X_s}c^n(X_s,\Omega_{X_s}^1)
=(-1)^{n+1}\chi(S)\chi(X_s).
\]

典范映射
\[
f^*:f^*\Omega_S^1\longrightarrow\Omega_X^1
\]
等同于 $\Omega_X^1\otimes(f^*\Omega_S^1)^{\otimes-1}$ 的一个截面 $a$. 在 $S$ 和 $X$ 上选择局部坐标系, 从而将此丛等同于 $\mathcal O_X^{n+1}$ 后, $a$ 的各坐标为偏导数 $\partial f/\partial x_i$. 因而 $a$ 只在 $F$ 上为零, 且由 1.3, 其在 $x\in F$ 处的零点重数就是 Milnor 数 $\mu(X/S,x)$.

回顾, 若纯维数为 $d$ 的光滑固有 $k$-概形 $X$ 上的秩 $d$ 向量丛 $G$ 的一个截面 $a$ 只有孤立零点, 则 $\int_Xc^d(G)$ 等于 $a$ 按重数计数的零点个数. 在这里,
\[
\mathrm{II}=\int_Xc^{n+1}\!\left(
\Omega_X^1\otimes(f^*\Omega_S^1)^{\otimes-1}\right)
=\sum_{x\in F}\mu(X/S,x).
\]
合并各项得
\[
\chi(X)=(-1)^{n+1}(\mathrm I+\mathrm{II})
=\chi(S)\chi(X_s)-(-1)^n\sum_{x\in F}\mu(X/S,x),
\]
由于 $\chi(X_s)=\chi(X_{\bar\eta})$, 这就证明了 2.1.
